\documentclass[../../main.tex]{subfiles}
\begin{document}
% SOURCE p=208; printed=193; object=ch06-inequalities-title; region=top
\section{几个有用的不等式}\label{section:ch06-useful-inequalities}
% SOURCE p=208; printed=193; object=ch06-inequalities-introduction; region=upper
估计与不等式是 $L^p$ 空间在分析中应用的核心，其中最基本的是 Hölder 与 Minkowski 不等式。本节给出这方面另外几个重要结果。第一个结果几乎显然，但十分有用，值得单独列出。
% SOURCE p=208; printed=193; object=ch06-chebyshev; number=6.17; region=upper
\begin{theorem}[Chebyshev 不等式]\label{theorem:ch06-chebyshev}
若 $f\in L^p$，其中 $0<p<\infty$，则对任意 $\alpha>0$，
\begin{align*}
\mu\{x:|f(x)|>\alpha\}\leqslant\left(\frac{\|f\|_p}{\alpha}\right)^p.
\end{align*}
\end{theorem}
% SOURCE p=208; printed=193; object=ch06-chebyshev-proof; region=middle
\begin{proof}
令 $E_\alpha=\{x:|f(x)|>\alpha\}$。则
\begin{align*}
\|f\|_p^p=\int|f|^p\geqslant\int_{E_\alpha}|f|^p
\geqslant\alpha^p\int_{E_\alpha}1=\alpha^p\mu(E_\alpha).
\end{align*}
除以 $\alpha^p$ 即得结论。
\end{proof}
% SUPPLEMENT: id=v14-ch06-chebyshev-conditions; parent=ch06-chebyshev; kind=proof-detail+motivation+hypotheses; origin=independent-derivation
\begin{remark}
\textbf{动机与逐条件说明（重构新增）。}一个函数若在很大的集合上超过固定高度，其 $p$ 次积分就不可能很小；此式把积分大小变为水平集测度的控制。可测性保证水平集可测；$p>0$ 保证取 $p$ 次幂保持大小关系，$p<\infty$ 使积分表达式适用；$\alpha>0$ 保证可除以 $\alpha^p$。$f\in L^p$ 使右端有限，因而每个正水平集测度有限；即使该积分无穷，扩展实数意义的不等式仍成立但不给有限估计。论证没有使用三角不等式，所以也包含 $0<p<1$，且不要求测度有限、半有限或 $\sigma$ 有限。把 $>$ 换为 $\geqslant$ 同样可证。
\end{remark}
% END SUPPLEMENT: id=v14-ch06-chebyshev-conditions
% SOURCE p=208; printed=193; object=ch06-kernel-introduction; region=middle
下面是关于 $L^p$ 空间上积分算子有界性的一个相当一般的定理。
% SOURCE p=208; printed=193; object=ch06-kernel-bound; number=6.18; region=middle
\begin{theorem}\label{theorem:ch06-kernel-bound}
设 $(X,\mathscr M,\mu)$ 与 $(Y,\mathscr N,\nu)$ 是 $\sigma$ 有限测度空间，$K$ 是 $X\times Y$ 上的 $(\mathscr M\otimes\mathscr N)$ 可测函数。假设存在 $C>0$，使
\begin{align*}
\int_X|K(x,y)|\,\mathrm d\mu(x)&\leqslant C\quad\text{对几乎所有 }y\in Y,\\
\int_Y|K(x,y)|\,\mathrm d\nu(y)&\leqslant C\quad\text{对几乎所有 }x\in X.
\end{align*}
若 $1\leqslant p\leqslant\infty$ 且 $f\in L^p(\nu)$，则
\begin{align*}
Tf(x)=\int_YK(x,y)f(y)\,\mathrm d\nu(y)
\end{align*}
对几乎所有 $x\in X$ 绝对收敛，由此定义的 $Tf$ 属于 $L^p(\mu)$，并且 $\|Tf\|_p\leqslant C\|f\|_p$。
\end{theorem}
% SOURCE p=208,209; printed=193,194; object=ch06-kernel-bound-proof; region=lower-top
\begin{proof}
先设 $1<p<\infty$，$q$ 是 $p$ 的共轭指数。对乘积
\begin{align*}
|K(x,y)f(y)|=|K(x,y)|^{1/q}\bigl(|K(x,y)|^{1/p}|f(y)|\bigr)
\end{align*}
应用 Hölder 不等式，得到对几乎所有 $x\in X$，
\begin{align*}
\int_Y|K(x,y)f(y)|\,\mathrm d\nu(y)
&\leqslant\left(\int_Y|K(x,y)|\,\mathrm d\nu(y)\right)^{1/q}
\left(\int_Y|K(x,y)||f(y)|^p\,\mathrm d\nu(y)\right)^{1/p}\\
&\leqslant C^{1/q}\left(\int_Y|K(x,y)||f(y)|^p\,\mathrm d\nu(y)\right)^{1/p}.
\end{align*}
所以由 Tonelli 定理，
\begin{align*}
\int_X\left(\int_Y|K(x,y)f(y)|\,\mathrm d\nu(y)\right)^p\mathrm d\mu(x)
&\leqslant C^{p/q}\int_X\int_Y|K(x,y)||f(y)|^p\,\mathrm d\nu(y)\,\mathrm d\mu(x)\\
&\leqslant C^{p/q+1}\int_Y|f(y)|^p\,\mathrm d\nu(y).
\end{align*}
最后的积分有限，故由 Fubini 定理可知 $K(x,\cdot)f\in L^1(\nu)$ 对几乎所有 $x$ 成立，于是 $Tf$ 几乎处处有定义，而且
\begin{align*}
\int_X|Tf(x)|^p\,\mathrm d\mu(x)\leqslant C^{p/q+1}\|f\|_p^p.
\end{align*}
取 $p$ 次方根即得结论。$p=1$ 时证明类似而更容易，只需假设 $\int_X|K(x,y)|\,\mathrm d\mu(x)\leqslant C$；$p=\infty$ 时只需假设 $\int_Y|K(x,y)|\,\mathrm d\nu(y)\leqslant C$，证明也很直接。原书将细节留作\cref{exercise:ch06-kernel-endpoints}。
\end{proof}
% SUPPLEMENT: id=v14-ch06-kernel-bound-completion; parent=ch06-kernel-bound; kind=proof-detail+motivation+hypotheses; origin=independent-derivation
\begin{remark}
\textbf{绝对收敛、良定义及逐条件说明（重构新增）。}令 $H(x)=\int_Y|K(x,y)f(y)|\,\mathrm d\nu(y)$。上面的计算先对非负函数用 Tonelli，允许积分为无穷，最后才得到 $\int H^p<\infty$，从而 $H<\infty$ 几乎处处。这正是原文援引 Fubini 处所需的结论；无需也不应额外声称 $K(x,y)f(y)$ 在整个乘积空间可积。由 $p/q+1=p$，最后的常数开方为 $C$。积分的实部与虚部各自正负部分可测，故在 $H<\infty$ 上得到可测的 $Tf$；在例外零集上令它为零。

两个端点的完整证明见本节习题解答。若 $f$ 在一个 $\nu$ 零集上改变，对每个可测截面的积分均不变；因此算子在 $L^p$ 等价类上良定义，线性来自绝对收敛积分的线性。联合可测性用于截面可测及 Tonelli；$\sigma$ 有限性是此处乘积积分定理的适用条件。列积分界控制 $L^1$，行积分界控制 $L^\infty$，两者经 Hölder 配合控制中间指数。$1\leqslant p\leqslant\infty$ 与 Hölder 的共轭指数相配；本证明没有给出 $p<1$ 的结论。$C$ 是统一的有限界，否则不能得出统一算子范数；$C>0$ 便于表达，$C=0$ 时同样由非负积分推出零算子。
\end{remark}
% END SUPPLEMENT: id=v14-ch06-kernel-bound-completion
% SOURCE p=209; printed=194; object=ch06-integral-minkowski-introduction; region=upper
Minkowski 不等式说，和的 $L^p$ 范数不超过各项 $L^p$ 范数之和。把求和换成积分，就得到下面的推广。
% SOURCE p=209; printed=194; object=ch06-integral-minkowski; number=6.19; region=middle
\begin{theorem}[积分形式的 Minkowski 不等式]\label{theorem:ch06-integral-minkowski}
设 $(X,\mathscr M,\mu)$、$(Y,\mathscr N,\nu)$ 是 $\sigma$ 有限测度空间，$f$ 是 $X\times Y$ 上的 $(\mathscr M\otimes\mathscr N)$ 可测函数。
\begin{enumerate}
\item\label{item:ch06-integral-minkowski-positive} 若 $f\geqslant0$ 且 $1\leqslant p<\infty$，则
\begin{align*}
\left[\int_X\left(\int_Y f(x,y)\,\mathrm d\nu(y)\right)^p\mathrm d\mu(x)\right]^{1/p}
\leqslant\int_Y\left[\int_X f(x,y)^p\,\mathrm d\mu(x)\right]^{1/p}\mathrm d\nu(y).
\end{align*}
\item\label{item:ch06-integral-minkowski-complex} 若 $1\leqslant p\leqslant\infty$，对几乎所有 $y$ 有 $f(\cdot,y)\in L^p(\mu)$，且函数 $y\mapsto\|f(\cdot,y)\|_p$ 属于 $L^1(\nu)$，则对几乎所有 $x$ 有 $f(x,\cdot)\in L^1(\nu)$，函数 $x\mapsto\int_Yf(x,y)\,\mathrm d\nu(y)$ 属于 $L^p(\mu)$，且
\begin{align*}
\left\|\int_Yf(\cdot,y)\,\mathrm d\nu(y)\right\|_p
\leqslant\int_Y\|f(\cdot,y)\|_p\,\mathrm d\nu(y).
\end{align*}
\end{enumerate}
\end{theorem}
% SOURCE p=209; printed=194; object=ch06-integral-minkowski-proof; region=lower
\begin{proof}
$p=1$ 时，(a) 就是 Tonelli 定理。若 $1<p<\infty$，令 $q$ 为其共轭指数，取 $g\in L^q(\mu)$。由 Tonelli 定理和 Hölder 不等式，
\begin{align*}
\int_X\left(\int_Yf(x,y)\,\mathrm d\nu(y)\right)|g(x)|\,\mathrm d\mu(x)
&=\int_Y\int_X f(x,y)|g(x)|\,\mathrm d\mu(x)\,\mathrm d\nu(y)\\
&\leqslant\|g\|_q\int_Y\left(\int_X f(x,y)^p\,\mathrm d\mu(x)\right)^{1/p}\mathrm d\nu(y).
\end{align*}
因此 (a) 由\cref{theorem:ch06-dual-test-criterion}得到。$p<\infty$ 时，对 $|f|$ 应用 (a)，再用 Fubini 定理即可得到 (b)；$p=\infty$ 时，(b) 是积分单调性的直接结果。
\end{proof}
% SUPPLEMENT: id=v14-ch06-integral-minkowski-completion; parent=ch06-integral-minkowski; kind=proof-detail+motivation+hypotheses; origin=independent-derivation
\begin{remark}
\textbf{对偶判别与无穷端点的完整补证（重构新增）。}在 (a) 中，右端为 $+\infty$ 时结论按扩展实数理解，已经成立。设右端 $A<\infty$，并令 $H(x)=\int_Y f(x,y)\,\mathrm d\nu(y)$。$H$ 是非负可测扩展值函数。$1<p<\infty$ 时，原证明的估计对每个有限测度集合 $E$ 的示性函数给出 $\int_EH\leqslant A\mu(E)^{1/q}<\infty$。用可数个这样的集合覆盖 $X$，可知 $H<\infty$ 几乎处处。在其无穷零集上改为零后，$H$ 满足对偶判别中所有有限支集简单测试的可积性及上界，故 $H\in L^p$ 且 $\|H\|_p\leqslant A$。这补上了引用判别前有限值的要求，也避免预先假设待证结论。

在 (b) 的有限 $p$ 情形，对 $|f|$ 使用 (a)，得到 $H(x)=\int_Y|f(x,y)|\,\mathrm d\nu(y)\in L^p$。所以截面积分绝对收敛几乎处处，且其模不超过 $H$，所求估计随即成立。可测性由正负实部和虚部的非负截面积分得到；无需假设 $f$ 在整个乘积空间属于 $L^1$。

现在设 $p=\infty$。选 $h(y)=\|f(\cdot,y)\|_\infty$ 的非负可测代表；依假设 $h\in L^1(\nu)$。对几乎所有 $y$，集合 $\{x:|f(x,y)|>h(y)\}$ 为 $\mu$ 零集。对联合可测坏集的示性函数使用 Tonelli，得到对几乎所有 $x$，$|f(x,y)|\leqslant h(y)$ 对几乎所有 $y$ 成立。因此
\begin{align*}
\int_Y|f(x,y)|\,\mathrm d\nu(y)\leqslant\int_Yh\,\mathrm d\nu<\infty
\end{align*}
对几乎所有 $x$ 成立，既给绝对收敛，也给所需无穷范数界。此处不可直接合并随 $y$ 变化的不可数多个 $\mu$ 零集。
\end{remark}
% END SUPPLEMENT: id=v14-ch06-integral-minkowski-completion
% SUPPLEMENT: id=v14-ch06-integral-minkowski-conditions; parent=ch06-integral-minkowski; kind=motivation+hypotheses; origin=independent-derivation
\begin{remark}
\textbf{动机与逐条件说明（重构新增）。}把一个积分看成连续叠加，范数的三角不等式就应成为“先取范数再积分”的上界；证明通过标量对偶测试落实这一想法，不预设向量积分理论。(a) 的非负性允许 Tonelli 与无穷积分；没有它又没有绝对可积性时，内积分可能根本不存在。$1\leqslant p$ 是范数三角不等式及 Hölder 估计所对应的范围；(a) 的有限 $p$ 使积分幂表达式有意义，无穷端点由 (b) 另行处理。(b) 的截面 $L^p$ 条件定义各个范数，范数函数属于 $L^1$ 则保证它可测且总量有限，这是绝对收敛的实际控制量。联合可测性与 $\sigma$ 有限性共同使交换非负积分和交换几乎处处量词合法。上述证明也适用于完成的乘积测度，改用\cref{theorem:ch02-complete-fubini-tonelli}即可；此时所有截面结论按几乎处处理解。
\end{remark}
% END SUPPLEMENT: id=v14-ch06-integral-minkowski-conditions
% SOURCE p=209; printed=194; object=ch06-homogeneous-kernel-introduction; region=bottom
最后一个结果讨论配备 Lebesgue 测度的 $(0,\infty)$ 上的积分算子。
% SOURCE p=210; printed=195; object=ch06-homogeneous-kernel; number=6.20; region=upper
\begin{theorem}\label{theorem:ch06-homogeneous-kernel}
设 $K$ 是 $(0,\infty)\times(0,\infty)$ 上的 Lebesgue 可测函数，对所有 $\lambda>0$ 满足
\begin{align*}
K(\lambda x,\lambda y)=\lambda^{-1}K(x,y).
\end{align*}
若对某个 $p\in[1,\infty]$，
\begin{align*}
\int_0^\infty|K(x,1)|x^{-1/p}\,\mathrm dx=C<\infty,
\end{align*}
并令 $q$ 为 $p$ 的共轭指数，则对 $f\in L^p$ 与 $g\in L^q$，令
\begin{align*}
Tf(y)=\int_0^\infty K(x,y)f(x)\,\mathrm dx,
\qquad Sg(x)=\int_0^\infty K(x,y)g(y)\,\mathrm dy.
\end{align*}
那么 $Tf$、$Sg$ 几乎处处有定义，且 $\|Tf\|_p\leqslant C\|f\|_p$、$\|Sg\|_q\leqslant C\|g\|_q$。
\end{theorem}
% SOURCE p=210; printed=195; object=ch06-homogeneous-kernel-proof; region=middle
\begin{proof}
令 $z=x/y$，则
\begin{align*}
\int_0^\infty|K(x,y)f(x)|\,\mathrm dx
&=\int_0^\infty|K(yz,y)f(yz)|y\,\mathrm dz\\
&=\int_0^\infty|K(z,1)f_z(y)|\,\mathrm dz,
\end{align*}
其中 $f_z(y)=f(yz)$；而且
\begin{align*}
\|f_z\|_p
=\left(\int_0^\infty|f(yz)|^p\,\mathrm dy\right)^{1/p}
=\left(\int_0^\infty|f(x)|^pz^{-1}\,\mathrm dx\right)^{1/p}
=z^{-1/p}\|f\|_p.
\end{align*}
故由积分形式的 Minkowski 不等式，$Tf$ 几乎处处存在，并且
\begin{align*}
\|Tf\|_p\leqslant\int_0^\infty|K(z,1)|\|f_z\|_p\,\mathrm dz
=\|f\|_p\int_0^\infty|K(z,1)|z^{-1/p}\,\mathrm dz=C\|f\|_p.
\end{align*}
最后令 $u=y^{-1}$，得到
\begin{align*}
\int_0^\infty|K(1,y)|y^{-1/q}\,\mathrm dy
&=\int_0^\infty|K(y^{-1},1)|y^{-1-1/q}\,\mathrm dy\\
&=\int_0^\infty|K(u,1)|u^{-1/p}\,\mathrm du=C.
\end{align*}
对转置核应用相同的论证，便知 $Sg$ 几乎处处有定义且 $\|Sg\|_q\leqslant C\|g\|_q$。
\end{proof}
% SUPPLEMENT: id=v14-ch06-homogeneous-kernel-completion; parent=ch06-homogeneous-kernel; kind=proof-detail+motivation+hypotheses; origin=independent-derivation
\begin{remark}
\textbf{端点与 Lebesgue 可测性的补全（重构新增）。}原证明中用 $p$ 次积分计算范数的等式适用于 $p<\infty$；若 $p=\infty$，伸缩 $y\mapsto zy$ 保持零测集及其逆像为零，故直接有 $\|f_z\|_\infty=\|f\|_\infty$，仍等于约定 $1/\infty=0$ 下的 $z^{-1/p}\|f\|_p$。对 $K(z,1)f_z(y)$ 应用 Minkowski 定理的 (b) 部分，所得首先是绝对收敛几乎处处，故符号积分也存在。转置计算中 $-1+1/q=-1/p$，包括 $p=1$ 或 $\infty$；不用把无穷指数代入有限幂范数公式。

还有一个完成测度下的细节：一般的平面 Lebesgue 可测函数不能不加说明地取任意固定截面。这里齐次性足以补上。映射 $(z,y)\mapsto(yz,y)$ 是正象限的 $C^1$ 微分同胚，在可数个紧矩形附近局部双 Lipschitz，因此它及逆映射保持零集；于是 $H(z,y)=K(yz,y)$ 是 Lebesgue 可测函数。由完成形式的 Fubini 定理，存在某个 $y_0>0$，使整个正半轴上的截面 $H(\cdot,y_0)$ 可测。齐次性给
\begin{align*}
k(z):=K(z,1)=y_0H(z,y_0),\qquad K(x,y)=y^{-1}k(x/y),
\end{align*}
故 $k$ 可测。用可数个矩形覆盖即可把截面的局部断言合为同一个 $y_0$。同样的坐标变换说明 $f(yz)$ 联合 Lebesgue 可测；或者先把 $f,k$ 换成几乎处处相等的 Borel 代表，再用完成形式的积分定理。每个固定正伸缩都保持一维零集，故改换这些代表不影响所需截面积分及其几乎处处类。
\end{remark}
% END SUPPLEMENT: id=v14-ch06-homogeneous-kernel-completion
% SUPPLEMENT: id=v14-ch06-homogeneous-kernel-conditions; parent=ch06-homogeneous-kernel; kind=motivation+hypotheses; origin=independent-derivation
\begin{remark}
\textbf{动机与逐条件说明（重构新增）。}负一次齐次性使积分核的缩放因子恰好与 $\mathrm dx=y\,\mathrm dz$ 抵消，于是算子变为伸缩函数 $f_z$ 的加权叠加。权 $z^{-1/p}$ 来自 Lebesgue 测度的伸缩规律，不能随意改掉；加权绝对积分有限正是 Minkowski 估计所需的可积控制。正半轴让比值、倒数和正伸缩均留在定义域；Lebesgue 测度使这些换元成立。$p\geqslant1$ 用于范数不等式，$q$ 共轭保证转置核的权变换仍给同一常数。齐次性是这里的充分结构条件，并不声称所有有界积分算子都必须齐次。$C=0$ 时估计仍成立；它使剖面为零几乎处处，算子为零。复数值核和函数也允许，因为控制量使用绝对值。
\end{remark}
% END SUPPLEMENT: id=v14-ch06-homogeneous-kernel-conditions
% SOURCE p=210; printed=195; object=ch06-hardy-inequalities; number=6.21; region=lower
\begin{corollary}\label{corollary:ch06-hardy-inequalities}
令
\begin{align*}
Tf(y)=y^{-1}\int_0^y f(x)\,\mathrm dx,\qquad
Sg(x)=\int_x^\infty y^{-1}g(y)\,\mathrm dy.
\end{align*}
那么当 $1<p\leqslant\infty$ 及 $1\leqslant q<\infty$ 时，
\begin{align*}
\|Tf\|_p\leqslant\frac{p}{p-1}\|f\|_p,\qquad
\|Sg\|_q\leqslant q\|g\|_q.
\end{align*}
\end{corollary}
% SOURCE p=210; printed=195; object=ch06-hardy-inequalities-proof; region=bottom
\begin{proof}
令 $K(x,y)=y^{-1}\chi_E(x,y)$，其中 $E=\{(x,y):x<y\}$。则
\begin{align*}
\int_0^\infty|K(x,1)|x^{-1/p}\,\mathrm dx
=\int_0^1x^{-1/p}\,\mathrm dx=\frac{p}{p-1}=q,
\end{align*}
这里 $q$ 是 $p$ 的共轭指数。由\cref{theorem:ch06-homogeneous-kernel}即得所求。
\end{proof}
% SUPPLEMENT: id=v14-ch06-hardy-conditions; parent=ch06-hardy-inequalities; kind=proof-detail+motivation+hypotheses; origin=independent-derivation
\begin{remark}
\textbf{动机、端点与条件说明（重构新增）。}$T$ 取从零到当前位置的平均，$S$ 取带倒数权的尾积分；它们由同一个三角区域上的核产生。核可测且为负一次齐次。$p>1$ 恰使剖面在零附近可积；$p=\infty$ 时应把上式解释为 $\int_0^1 1\,\mathrm dx=1$，常数为一，而不是运算 $\infty/(\infty-1)$。对 $T$ 的各个指数分别配上共轭 $q$，就遍历了 $S$ 的全部 $1\leqslant q<\infty$，并不要求陈述中独立给定的两次应用使用同一对函数。遗漏的端点一般不成立：$f=\chi_{(0,1)}\in L^1$ 时，$Tf(y)=1/y$ 在 $y>1$ 上不可积；$g=1\in L^\infty$ 时，$Sg(x)=\infty$ 对每个 $x>0$ 成立。对所述合法指数，绝对收敛和可测性由前一定理同时保证。
\end{remark}
% END SUPPLEMENT: id=v14-ch06-hardy-conditions
% SOURCE p=211; printed=196; object=ch06-hardy-name; region=top
\cref{corollary:ch06-hardy-inequalities}是 \textbf{Hardy 不等式}的一个特例；一般结果见\cref{exercise:ch06-weighted-hardy}。
\subfile{section-03-exercises.tex}
\end{document}
